(a) The polynomial $-(x_1^2+\cdots+x_r^2)$ is square-free for $r\geq2$: a repeated irreducible factor would divide all its partial derivatives $2x_1,\ldots,2x_r$. Thus \exref{II.6.4} applies.

(b) For (b)--(d), assume $k$ is algebraically closed, as in (I, Ex. 5.12). Choose $\iota\in k$ with $\iota^2=-1$. Replacing $x_0,x_1$ by $x_0+\iota x_1,-x_0+\iota x_1$ gives $x_0x_1=x_2^2+\cdots+x_r^2$. The unused coordinates contribute affine-space factors and do not change the class group, by (6.6). Write $A$ for the resulting coordinate ring. Since $A_{x_0}=k[x_0,x_0^{-1},x_2,\ldots,x_n]$ is a UFD, (6.5) says that $\operatorname{Cl}X$ is generated by the components of $x_0=0$.

For $r=2$, there is one component $D=V(x_0,x_2)$, and $(x_0)=2D$. Its ideal is not principal: its image in $\mathfrak m/\mathfrak m^2$ at the origin contains the independent classes of $x_0,x_2$, whereas a principal ideal contained in $\mathfrak m$ has image of dimension at most one. By normality and (6.2), $D$ is not principal. Hence $\operatorname{Cl}X\cong\mathbb Z/2\mathbb Z$.

For $r=3$, factor $x_2^2+x_3^2=uv$. The two components $D=V(x_0,u)$ and $E=V(x_0,v)$ satisfy $(x_0)=D+E$. If $aD+bE=(f)$, then $f$ is a unit on $D(x_0)$, so $f=cx_0^m$ for $c\in k^*$ and $m\in\mathbb Z$. Thus $a=b=m$, and $\operatorname{Cl}X\cong\mathbb Z^2/\mathbb Z(1,1)\cong\mathbb Z$.

For $r\geq4$, $x_2^2+\cdots+x_r^2$ is irreducible, since a reducible quadratic has rank at most two. Thus $x_0=0$ is a single reduced prime divisor, equal to $(x_0)$, and $\operatorname{Cl}X=0$.

(c) First take $n=r$. For $r=2$, the conic is the degree-two image of $\mathbb P^1$, so $\operatorname{Cl}Q\cong\mathbb Z$ and $[Q.H]$ is twice the generator. For $r=3$, the Segre isomorphism gives $Q\cong\mathbb P^1\times\mathbb P^1$, so $\operatorname{Cl}Q\cong\mathbb Z^2$ by (6.6.1). For $r\geq4$, \exref{II.6.3}(b) and part (b) give $\operatorname{Cl}Q=\mathbb Z[Q.H]$. If $n>r$, adding each unused homogeneous coordinate forms a projective cone. By \exref{II.6.3}(a),(b), this preserves both the class group and its hyperplane class, giving the same conclusions.

(d) By (a),(b) and (6.2), $S(Q)$ is a UFD for $r\geq4$. The homogeneous height-one prime defining $Y$ is therefore $(g)$ for a homogeneous element $g$. Lift $g$ to a homogeneous polynomial $G$ on $\mathbb P^n$. If $G$ factored into two positive-degree homogeneous factors, their images would factor $g$ into nonunits, contradicting primality. Thus $G$ is irreducible. The ideal of the scheme-theoretic intersection $V_+(G)\cap Q$ in $S(Q)$ is $(g)$, so this intersection is $Y$, with multiplicity one.

Without the algebraically closed-field hypothesis, the coordinate change in (b) need not exist: over $\mathbb R$, a sum of squares has no nonzero zero, whereas $x_0x_1=x_2^2+\cdots+x_r^2$ does.
